\documentclass[11pt,a4paper]{article}
\usepackage{amsmath,amssymb,amsthm}
\usepackage[margin=2.5cm]{geometry}
\usepackage{hyperref}

\newtheorem{definition}{Definition}[section]
\newtheorem{theorem}[definition]{Theorem}
\newtheorem{proposition}[definition]{Proposition}
\newtheorem{remark}[definition]{Remark}
\newtheorem{conjecture}[definition]{Conjecture}

\title{Hyperseed Formalization:\\
  The General Theory of General Intelligence}
\author{ProtomegaTron (automated formalization)}
\date{2026-09-18}

\begin{document}
\maketitle

\begin{abstract}
We formalize Ben Goertzel's ``The General Theory of General Intelligence'' (2021-06-02)
within the Hyperseed ontology.
\end{abstract}

\section{Intelligence as fiber richness}

\begin{proposition}[Intelligence = thick fiber --- claims 1, 4, 7]
General intelligence is thick, diverse fiber over base space:
\[
  \text{GI}(b) = \{F_i(b)\}_{i \in I} \quad \text{with} \quad |I| \gg 1
\]
The richness of fiber (number and diversity of fibers over each base point)
determines the system's ability to respond flexibly to diverse environmental
conditions. Thin fiber = narrow intelligence; thick fiber = general intelligence.
\end{proposition}

\section{Functors as fiber maps}

\begin{proposition}[Functors = fiber transport --- claims 2, 8]
Category-theoretic functors between environment and action categories
correspond to coherent fiber transport:
\[
  \mathcal{F}: \text{Env} \to \text{Act} \quad \leftrightarrow \quad
  \text{parallel transport along } \gamma \in B
\]
Natural transformations between functors correspond to learning ---
changing the fiber map while preserving coherence.
\end{proposition}

\section{Self-modification as fiber evolution}

\begin{proposition}[Self-modification = autopoietic fiber --- claims 3, 9]
Self-modification is the fiber reshaping itself:
\[
  \frac{d}{dt} F(b,t) = \Phi(F(b,t), b) \quad \text{(autopoietic dynamics)}
\]
The evolution operator $\Phi$ depends on the current fiber state ---
the system modifies its own structure based on its current structure.
\end{proposition}
\end{document}
